\chapter{Spatiotemporal mobility patterns in \textit{Bicing}} \label{ch:bicing_mobility_patterns}

\textcolor{red}{
    This Chapter provides the mobility context for the battery and \ac{pm} frameworks developed in Chapters~\ref{ch:mobility_and_battery_modeling_bicing} and~\ref{ch:failure_forecasting_bicing}.
}

\textcolor{blue}{
    Part~\ref{pt:predictive_operations} turns to developing operational predictions for \textit{Bicing}: bicycle flows, \ac{e-b} batteries, and \ac{pm} for bike components. After the case-study introduction to the system, this Chapter provides the mobility context needed for the predictive frameworks developed in Chapters~\ref{ch:mobility_and_battery_modeling_bicing} and~\ref{ch:failure_forecasting_bicing}.
}
The material in this Chapter is adapted from Grau-Escolano et al.~\cite{Grau-Escolano2024}, published in \textit{EPJ Data Science}.

\section{Introduction}

\acp{bss} have been extensively studied from various angles, including station location optimization~\cite{Liu2015} and fairness~\cite{DuranRodas2020}, bike rebalancing strategies~\cite{Chardon2016}, as well as studies on user behavior changes and their impact on health~\cite{Ma2020, Qiu2018}. While substantial research has focused on human mobility within \acp{bss} in recent years~\cite{Talavera_Garcia2021, Kon2021}, the advent of fourth-generation \acp{bss} has introduced a new layer of complexity through the global adoption of \acp{e-b}~\cite{Shaheen2010, MeddinWorldMapReport2022}. Existing \acp{bss}, such as Barcelona's \textit{Bicing}, in Spain, have progressively integrated \acp{e-b} into their \acp{m-b} fleets, while there are instances of entirely new fully electric systems like Barcelona metropolitan area's \textit{AMBici}. As \ac{bss} users increasingly embrace \acp{e-b} for their speed, convenience, and reduced physical exertion, significant changes in existing \ac{bss} mobility dynamics are expected. However, the evolving mobility dynamics resulting from the coexistence of both mechanical and electric transportation modes remain an area ripe for study.

This Chapter addresses that gap by characterizing the main spatiotemporal mobility patterns of \textit{Bicing} during 2022, with a particular focus on how \acp{m-b} and \acp{e-b} differ in terms of trip counts, duration, distance, speed, and elevation, and how topography shapes the spatial distribution of both modes across stations. \textcolor{red}{
    The analysis provides the empirical context needed for the battery and \ac{pm} frameworks developed in Chapters~\ref{ch:mobility_and_battery_modeling_bicing} and~\ref{ch:failure_forecasting_bicing}.
}

\section{Literature review}

Research in \ac{bss} mobility typically involves the examination of aspects like temporal usage patterns and trip characteristics, which include factors such as distance, duration, and speed. These kinds of analyses serve multiple purposes, such as examining general system dynamics~\cite{Jensen2010, Ciancia2015, Zaltz2013}, designing effective rebalancing strategies~\cite{Chiariotti2018}, predicting \ac{bss} demand~\cite{Borgnat2011}, and estimating trip destinations and durations~\cite{Zhang2016Bicycle}. While trip characteristics can vary depending on the specific city in which the \ac{bss} is implemented, existing publications generally concur on certain approximate values. On average, \ac{bss} trips tend to have a mean distance of between 1 and 2 kilometers, a mean duration ranging from 10 to 20 minutes, and an average speed in the range of 10 to 15 kilometers per hour. Regarding trip distances, \ac{gps} data has been used to map actual ride routes~\cite{Talavera_Garcia2021}. However, when \ac{gps} data is unavailable, trip distances are estimated. In some instances, the shortest paths between origin and destination stations are calculated~\cite{Oliveira2016}. Nevertheless, comparisons between \ac{gps} routes and the shortest paths have suggested that the shortest paths may not accurately reflect the actual choices made by walkers, drivers, and cyclists~\cite{Zhu2015, Lu2018, Klein2023}. Factors such as a greener environment, the presence of amenities, or roads with greater connectivity often lead to deviations from the shortest paths~\cite{Klein2023}.

An alternative perspective for studying \ac{bss} mobility emphasizes the geographical aspects of mobility. While some studies have examined the overall usage of \acp{bss} across multiple cities~\cite{Zhao2014, DuranRodas2019, Sarkar2015}, much of this research is centered around understanding the patterns within individual \acp{bss}. Most previous studies in this domain focus on the usage of individual docking stations, often utilizing concepts such as incoming and outgoing trips~\cite{Moncayo2016, FaghihImani2014, Wang2016, FaghihImani2017Empirical, DuranRodas2019, Oliveira2016}, as well as station occupancy~\cite{OBrien2014, Sarkar2015}. However, for a more comprehensive understanding of \ac{bss} mobility structures, the flows between stations have also been explored~\cite{Corcoran2014, Yang2019, Kon2021}.

A relatively less studied area is the relationship between urban topology and \ac{bss} usage, specifically how the altitude difference between the \ac{od} stations of a trip influences bike usage. In this regard,~\cite{Morency2015} found that elevation has a negative impact on the number of incoming trips and a positive impact on the outgoing, while~\cite{Kim2020} found that altitude difference, together with stations' distance and weather features are good predictors of bike usage.

On the other side, when characterizing \ac{bss} mobility, the majority of research has predominantly focused on traditional \acp{m-b}, with only a limited number of studies comparing their usage with other transportation modes. For instance,~\cite{Noussan2019} analyzed the temporal variations in \acp{bss} and car usage, while~\cite{Jensen2010} compared \ac{bss} trip distances and average speeds with those of car and pedestrians. Moreover, there is an important field of investigations in comparing between \acp{e-b} and scooters. Spatio-temporal mobility patterns of \acp{bss} and dockless scooter-sharing services were compared in~\cite{Mckenzie2019}, finding that \acp{bss} are mainly used for commuting, whereas scooter-sharing serves different purposes. Additionally,~\cite{Almannaa2020} explored the average speed of shared e-scooters and \ac{bss} \acp{e-b}, finding that \acp{e-b} generally travel faster than e-scooters.

Mobility data from \textit{Bicing} has served as a valuable resource in various articles, with the majority of studies relying on data from around 2010, a period characterized by fewer stations, bikes, and bike lanes compared to the situation in 2022. Notably, this was also a time before the introduction of \acp{e-b} into the \textit{Bicing} system. In this context, two early studies obtained from the \textit{Bicing} website the number of occupied and empty bike docks at stations with the purpose of unraveling the general spatio-temporal patterns of Barcelona dynamics~\cite{Froehlich2008} and comparing \ac{bss} usage patterns on weekdays and weekends employing hierarchical clustering~\cite{Froehlich2009}. Additionally, Bayesian networks were used to predict the number of available bicycles at each station. In the same line,~\cite{Kaltenbrunner2010} utilized the same data to generate bike availability predictions in the stations using \ac{arma} models. More recently, in 2022, a study delved into predicting the usage of the \ac{bss} system and examined the impact of the COVID-19 pandemic on these predictions~\cite{Bustamante2022}. This last research used a much more comprehensive and updated dataset, including data from 2020 and 2021, which contained information about the \ac{od} stations, as well as the start and end times of individual trips.



\section{Methodology}

\subsection{Data}

The analysis is based on trip records provided by \textit{Bicing}, which encompass the following information:

\vspace{-1.1em}
\begin{itemize}\setlength\itemsep{0em}
    \item Starting and ending dates and times with second-level granularity
    \item Starting and ending stations
    \item Bike identifier
    \item Bike model (\ac{m-b} or \ac{e-b})
    \item Anonymized user identifier
\end{itemize}
\vspace{-0.6em}

The complete dataset covers trips from April 2019 to December 2022 (i.e., 3 years and 9 months), totaling 53 million trips (62\% \ac{m-b}, 38\% \ac{e-b}). However, this Section uses only trips from 2022 to minimize the impact of the COVID-19 pandemic and to reflect the increasing use of \acp{e-b}, as detailed in Appendix~\ref{app_paper3:section_data_exploration}. Additionally, geographical data for the 519 stations, including latitude and longitude coordinates, are provided.

    
    \subsection{Trips processing} \label{section_trips_processing}
  
    To study the mobility patterns, only trips with durations ranging from 2 to 60 minutes were taken into account, which constituted 99\% of all trips. Once filtered, various trip metrics were collected for the mobility analysis. While trip duration could be directly derived from the trips data, other measures required some additional steps. Obtaining the trip routes was not feasible since data only provided the starting and ending stations for each trip.

    As a consequence, the distances between stations were obtained using the OpenRouteService \ac{api} \cite{OpenRouteService_2023}, which can suggest recommended routes for bicycles based on \ac{osm} data~\cite{OpenStreetMap_2023}. These routes are generated by combining the suitability of streets for the chosen mode of transportation along with the fastest route option. 

    Then, using the trips' duration and distance their average speed was obtained. Furthermore, the cumulative trip inclinations were simplified by calculating the difference between the altitudes of the origin and destination stations. The specific station altitudes were obtained using the \ac{otd} \ac{api}~\cite{OpenTopoData_2023}.

    To determine whether the trip characteristics of \acp{m-b} and \acp{e-b} come from the same distribution, 5,000 samples from each subgroup were randomly selected. Initially, a Shapiro-Wilk test was applied to both samples to assess whether they followed a normal distribution. Subsequently, if both samples were found to be normally distributed, an Independent Sample T-Test was applied for comparison. In cases where one or both samples did not meet the normality assumption, the Kolmogorov-Smirnov test was used. All statistical tests were conducted with a significance level set at 0.01.

    To compare trip numbers of both bike models, several metrics were computed to avoid comparisons with absolute numbers. Initially, for each bike model, at all stations, the percentage of incoming (or outgoing) trips was determined by dividing the total trips for each model by the station's total trips. Then, two supplementary metrics were derived from these percentages: (1) the differences in the percentage of incoming (or outgoing) trips between \acp{e-b} and \acp{m-b}, and (2) the difference in the percentage of incoming and outgoing trips for each bike model.


\section{Results}

Both mechanical and electric transportation modes exhibited variability across time due to seasonal variations and holidays such as Easter week and summer holidays (Figure~\ref{fig:trips_characteristics}A). However, a noticeable differential trend emerged from summer onwards. Electric mobility experienced a rise in trip numbers while the mechanical one suffered a significant decrease. Despite some \acp{m-b} being upgraded to \acp{e-b} in the course of 2022, the increase in the number of \ac{e-b} rides couldn't be entirely attributed to this transformation, considering they comprised only 40\% of the bike fleet at their highest point. The \acp{e-b} experienced a rise in the mean daily trips per bike, whereas the \acp{m-b} saw a significant reduction. Consequently, the variations in the usage of both transportation modes can be attributed to the upgrade of \acp{m-b} into \acp{e-b} and the increasing preference of users for the electric mobility.

\begin{figure}[ht!]
    \centering
    \includegraphics[width=1\textwidth]{0_plots/paper3/main/panel_1.png}
    \caption[Trip patterns of \textit{Bicing} in 2022]{
        \textbf{Trip patterns of \textit{Bicing} in 2022.} 
        \textbf{(A)} Time series depicting the weekly trip counts alongside the average weekly number of trips per bike. 
        \textbf{(B)} Hourly time series illustrating the average trip counts throughout the week, with the area around the time series indicating the standard deviation. 
        \textbf{(C-F)} Distributions of trip characteristics, with vertical lines representing the mean of each subgroup.
    }
    \label{fig:trips_characteristics}
\end{figure}    

Trip counts present a strong weekly seasonality characterized by two distinct mainly daily patterns: weekdays and weekends (Figure~\ref{fig:trips_characteristics}B). On weekdays, three important trip count maxima are observed at 8:00, 14:00, and 18:00, with a gradual increase in trip numbers as the week progresses, peaking on Thursday. Fridays slightly deviate from the typical weekday pattern since they present similar trip count peaks at 14:00 and 18:00. In contrast, weekends are characterized by a 30\% mobility reduction, the absence of an 8:00 peak, and a shift in the 18:00 peak to 19:00. Notably, electric mobility consistently matches or surpasses the mechanical one across all days and hours despite its additional cost. Both transportation modes also differ in their trip characteristics (Figure~\ref{fig:trips_characteristics}C-F). Significant differences were found in the distributions of trip duration, distance, speed, and elevation (Section~\ref{section_trips_processing}), implying two distinct behaviors. Specifically, electric mobility is characterized by steeper inclines, longer durations, greater distances from the origin, and higher speeds compared to mechanical mobility.

Additionally, each mobility mode exhibits unique spatial mobility dynamics (Figure~\ref{fig:elevation_maps}). High-altitude stations predominantly receive \ac{e-b} trips, whereas low-altitude stations see the majority of incoming trips via \acp{m-b}. This reflects users' preference for electric mobility when doing physically demanding trips. While this pattern is clear for incoming trips, it becomes less pronounced for outgoing trips, since the proportional utilization of \acp{e-b} is not as dominant, even at high-altitude stations. This is attributed to the fact that the outgoing trips from a station do not indicate their destinations, resulting in a mix of trips to lower and higher stations. Nevertheless, given the overall preference for electric mobility among users, the inclination towards \acp{e-b} still remains evident. Furthermore, when analyzing the differences in the percentages of incoming and outgoing trips involving \acp{e-b} (Appendix~\ref{app_paper3:section_mobility_analysis}), it becomes evident that the prevalence of incoming \ac{e-b} trips is more pronounced at high-altitude stations. In contrast, at lower-altitude stations, the percentages of incoming trips made by \acp{e-b} become less substantial compared to the outgoing trips.

\begin{figure}[ht!]
    \centering
    \includegraphics[width=1\textwidth]{0_plots/paper3/main/panel_2.pdf}
    \caption[Trip flow percentages for e-bikes and m-bikes in 2022, segmented by trip elevation.]{
        \textbf{Trip flow percentages for \acp{e-b} and \acp{m-b} in 2022, segmented by trip elevation.} Each pair of maps displays the differences in \ac{e-b} and \ac{m-b} trip percentages for incoming and outgoing trips at each station. The left maps include all trips, whereas the remaining figures apply a filter based on trip elevation. Within each map, individual dots represent \ac{bss} stations. The dot size reflects the station's altitude, with larger dots indicating higher altitudes, and the dot color signifies differences in trip flow percentages. Gray dots denote stations that do not receive or send trips with the specified altitude filter.
    }
    \label{fig:elevation_maps}
\end{figure}

By segregating incoming and outgoing trips based on their elevation, this phenomenon becomes even more evident (Figure~\ref{fig:elevation_maps}). At stations that generate outgoing trips with an elevation gain of over 100 meters, electric mobility completely dominates the transportation mode. Consequently, the receiving stations at higher elevations predominantly receive \ac{e-b} trips. This prevalence of electric mobility reduces gradually as the elevation decreases until reaching those trips with elevations between -50 and 50 meters. In this range of elevations, mechanical mobility plays a much more prominent role, especially at low altitude stations, however, even at high-altitude stations, electric mobility remains relevant in both incoming and outgoing trips. Notably, even in scenarios with negative elevations exceeding -50 meters, electric mobility continues to be the preferred transport mode choice, albeit with reduced percentages. This phenomenon may be attributed to users' preference for an electric transportation mode on uphill rides, resulting in more \acp{e-b} accumulating at higher altitudes, and to potential elevation irregularities in paths between high-altitude stations.

\section{Discussion}

The mobility analysis presented in this Chapter offers a novel perspective in \ac{bss} research by analyzing mobility patterns, differentiating between \acp{m-b} and \acp{e-b} as separate modes of transportation. Our findings reveal that \textit{Bicing} exhibits mobility dynamics similar to those in other \ac{bss} studies in terms of distance traveled, trip duration, and bike speed~\cite{Jensen2010, Ciancia2015, Zaltz2013, Chiariotti2018, Borgnat2011, Zhang2016Bicycle}. Notably, electric mobility has been found to be the preferred mode of transportation, characterized by longer trip durations, farther distances covered, and faster speeds when compared to mechanical mobility. Moreover, the trend of lower trip counts on weekends compared to weekdays, as noted in~\cite{Froehlich2009, Bustamante2022}, still holds true, along with the occurrence of three daily peaks~\cite{Froehlich2009}. 

Aligned with the findings of~\cite{Morency2015, Kim2020}, the analysis further confirms that topography plays a significant role in \ac{bss} utilization. Electric mobility is primarily chosen for trips involving steep positive inclines, whereas mechanical mobility is preferred for routes with slight elevation changes, especially for journeys originating or ending at stations situated at lower elevations. Consequently, stations at higher altitudes predominantly receive \acp{e-b}, whereas those at lower altitudes are more likely to attract the mechanical ones. Moreover, this pattern is even more pronounced when considering the elevation differences between the origin and ending stations. Interestingly, even in scenarios with significant negative elevations, where physical exertion is low, electric mobility remains the dominant choice. This preference may be attributed to the accumulation of \acp{e-b} at higher altitudes and the irregularities in the paths connecting high-altitude stations.

Despite the quality of the data, certain limitations may have affected this analysis. The absence of specific bike routes required us to infer trip distances (e.g., via routing APIs), potentially affecting the accuracy of mobility metrics such as distance and speed. Similarly, elevation gains were simplified to the difference between origin and destination station altitudes, potentially overlooking elevation irregularities along the paths between stations.

Future research could use the mobility insights from this Chapter to enhance bike rebalancing strategies for \acp{bss} that include both \acp{m-b} and \acp{e-b}, for example by optimizing bike distribution based on the detailed usage patterns documented here.